\begin{wrapfigure}[18]{r}{0.45\textwidth} 
    \vspace{-20pt} % Pulls the box up to align with the header
\begin{tcolorbox}[colback=sapienza!5, colframe=sapienza, title=Key Consideration, fonttitle=\bfseries]
\small
It is important to note that while the financial incentives and grants, the policimakers can’t just "throw money" at the problem and expect a perfect industrial symbiosis to appear. Subsidies might help buy the right technology, but they cannot buy trust between companies or fix the issue if factories are too far apart. For symbiosis to actually work, a complex mix of factors has to come together at the same time. The companies need to be physically close enough to exchange materials, their waste streams have to be technically compatible with their neighbors' needs, and the managers need to actually trust that their partners will be reliable. Financial policies are just the spark to get things moving; if those physical and social conditions aren't there, no amount of tax breaks or funding will create a functioning ecosystem.
\end{tcolorbox}
\end{wrapfigure}
Based on the previous analysis of the Ecodistretto di Porto Marghera, it is clear that relying only on factories and technology is not enough: the right policies are necessary to actually make the system run smoothly. To improve the current symbiosis between Veritas, Ecopatè, and the other industries in the area, several policy approaches that could help were identified. In the following paragraphs, some key areas where policy can make a difference are listed.
\subsection{EU Targets as a driver for cooperation}
The ambitious recycling targets set by the EU can be a driver for the existing IS of Porto Marghera. These regulations force companies to work together because with high targets for recovery single company usually can't manage it alone, they have to cooperate. For example, the waste management plants need to work closely with the glass, metal and plastic reclycling to reduce costs and increase the recovery percentage. Future policies should keep this pressure, effectively making "going it alone" impossible. But it also has to be taken into consideration that those polices can also be seen as barriers if certain systems are already implemented and a change of the processes requires a lot of resources and planning. 
\subsection{Making disposal expensive (Green Taxes)}
For industrial symbiosis to work, it has to be cheaper to reuse waste than to just throw it away. For example, turning waste into CSS (fuel) works because landfilling or exporting that waste would cost too much. Policy needs to keep disposal taxes high. If it becomes cheap to dump untreated waste in a landfill again, the economic motivation to exchange by-products disappears. Basically, high taxes on the "bad" options make the "symbiotic" options the only smart business choice.
\subsection{Variable fees based on quality}
A big problem for Porto Marghera is the quality of the materials. There is still a lot of contamination (like plastics mixed with glass) which makes recycling really hard and expensive (Gruppo Veritas, 2020, p.21). A good policy fix would be "variable waste fees." Instead of a flat rate, companies or municipalities should pay less if they deliver high-quality, clean waste, and pay more if the waste is contaminated. This would motivate better separation at the source, meaning facilities like Ecopatè get cleaner inputs that are easier to reuse.
\subsection{Long-term stability}
Companies are scared of uncertainty: if a company thinks the rules are going to change in two years, they will not be willing to invest millions in building a pipeline to a neighbor's factory. Policymakers need to give clear, long-term timelines so that firms feel safe investing in symbiosis for the long haul.
\subsection{Better information sharing}
Finally, policy should encourage or even require industries to share data on their waste flows openly, so that enterprises can exchange materials because they are aware of its quality. Veritas Group performs over 1,100 merceological analyses every year on separate waste collections (Gruppo Veritas, 2021, p.14), but this data needs to be more visible. If every company knows exactly what materials are available and what quality they are, it builds trust and makes it easier to spot new opportunities for cooperation.
\subsection{Subsidies for tech upgrades}
Upgrading the quality of the collected waste can be expensive. For example, buying new machinery and specific trucks requires a lot of money. Policies that offer grants or low-interest "green loans" can help overcome this. If the government helps reduce the risk of buying new technology, companies are much more likely to try to integrate their processes with their neighbors.
